% THETA3 contribution by /root/theta_032, 21 November 2026.
% GPT-6 Astra Pro; inherited reasoning-effort label not separately exposed.
% External theorems are cited inputs. A1--A5 deductions were checked in full.

\subsection{The exact theta dictionary and an applicable Rallis formula}

The attached 032 paper fixes a Galois CM field $E$, its real subfield
$F=E^+$ with $[F:\Q]\geq2$, and a Hermitian $E$-space $V$ of rank three,
of signature $(2,1)$ at the distinguished embedding $\mu$ and $(3,0)$
elsewhere. Its group is $\SU(V)$, and its domain $D$ parametrizes negative
lines $L\subset V_\mu$. Put $Q=V_\mu/L$; thus
\[
 T_D=\Hom(L,Q),\qquad \Omega_D^1=L\otimes Q^*.
\]
For a multiplicity space $P$ of rank $k$, the active positive Lagrangian
has pieces $L\otimes P_\mu$ and $Q^*\otimes P_\mu^*$. The half-determinant
line is $L^k$ times a constant line. These are the precise identifications
in 032, Section 4.1 and Lemma 4.1, printed pp.13--16. The character signs
are $d(\tau)=v(\tau^{-1})$, not $v(\tau)$. The finite Schwartz input, the
rational symplectic space, the embedding of $\SU(V)$, and $V$ itself remain
fixed in Proposition 4.6. Only multiplicity subspaces at compact places
move. This operation does not change the signature of $V$.

The first Taylor tensor projects to a holomorphic one-form. Proposition
4.5, pp.21--22, proves that finitely many actual finite inputs span any
specified cotangent fibre; the common level is chosen after selecting
those inputs. For $P=P_1\perp P_2$, double alternation of the completed
second Taylor tensor gives a two-form. Proposition 4.6, pp.22--24, proves
the product identity. In its divided-power convention the degree-two
coefficient of a product is
\[
 T_2^{(1)}T_0^{(2)}+T_1^{(1)}T_1^{(2)}+T_0^{(1)}T_2^{(2)}.
\]
Double alternation kills the two unmixed terms and sends the mixed term
to the wedge of the two first tensors with coefficient one. The general
completed second tensor need not be holomorphic. Intermediate compact
parameters are asserted only to give smooth forms; the rational
decomposition endpoints are holomorphic. Theorem 5.1, p.26, proved in
Section 6.3, and Corollary 6.7, p.45, then extract the rational CM Hodge
maps on a common principal cover. These arithmetic extraction results
are distinct from the analytic product identity.

The companion paper of van Hoften--Sempliner makes the global character
data explicit. In Sections 3.1 and 4.1, the $\mu_i$ are unitary
conjugate-symplectic Hecke characters of weight one, with restriction
$\epsilon_{E/F}$; their CM types satisfy
$\Phi_{\mu_i}=\Phi_i^{-1}$. The pair consists of $\U(V)$ and a
skew-Hermitian line $W_i$, with splitting characters $(\mu_i,\nu_5)$.
The character $\pi_i$ on $[\U(W_i)]$ is a source automorphic character,
and is not the splitting character $\mu_i$. Section 4.2 imposes the
global equality $\mu_1\mu_2=\mu_3\mu_4$. Lemma 4.2.1 supplies
$W_1\perp W_2\simeq W_3\perp W_4$. Doubling gives a split rank-four
skew-Hermitian space and the model $\mathcal S(V(\A_E)^2)$.
Remark 4.2.5 explains the cancellation of the diagonal $\U(V)$ character.
With that cancellation its action is $\Phi(x)\mapsto\Phi(g^{-1}x)$.
The doubling map has normalization
\[
 \delta(f\otimes\bar f)(0)=\|f\|_2^2.
\]
Proposition 4.2.6 and Section 4.2.8 express the fourfold cup period as a
specified geometric scalar times a doubled theta integral. Lemmas
4.2.7, 4.2.9, and 4.2.10 provide, respectively, nonvanishing, the
archimedean zero evaluation, and finite torus-character extraction.

For the rank-one to rank-three lift, an applicable explicit Rallis
formula is Haan, \emph{The Bessel period of $\U(3)$ and $\U(2)$ involving
a non-tempered representation}, arXiv:1403.5061v7, Section 3.3,
Theorem 3.4, printed p.8. Its proof cites J.-S. Li, \emph{Non-vanishing
theorems for the cohomology of certain arithmetic quotients}, Theorem
2.1, and the doubling factorization. This is a convergent-range
formula; no appeal to a second-term-range theorem with incompatible
dimensions is needed.

Here is its normalization. Let $V_1$ be Hermitian of rank one and $W_3$
skew-Hermitian of rank three, and put $T=\U(V_1)$ and $G=\U(W_3)$.
Haan uses the symplectic form
$\frac12\Tr_{E/F}(h_{V_1}\otimes h_{W_3})$, a fixed nontrivial additive
character $\psi$, and a unitary character $\gamma$ of
$E^\times\backslash\A_E^\times$ with
$\gamma|_{\A_F^\times}=\epsilon_{E/F}$. The splitting pair is
$(\gamma,\gamma^3)$. The theta lift is
\[
 \theta(\bar f,\varphi)(g)
 =\int_{[T]}\theta(t,g;\varphi)\overline{f(t)}\,dt,
\]
with Tamagawa measure. If $\tau$ is a unitary automorphic character of
$T$, $f_i=\bigotimes_v f_{i,v}\in\tau$, and
$\varphi_i=\bigotimes_v\varphi_{i,v}$, then
\begin{equation}\label{A-Rallis}
 \left\langle\theta(\bar f_1,\varphi_1),
                    \theta(\bar f_2,\varphi_2)\right\rangle_G
 =\frac{L_E(3/2,\operatorname{BC}(\tau)\otimes\gamma^3)}
        {L_F(3,\epsilon_{E/F})}
       \prod_v Z_v^{\#}(1,f_{1,v},f_{2,v},\Phi_{1,v}).
\end{equation}
The target Petersson measure is also Tamagawa. Local pairings are
chosen so that their product is the global Petersson pairing on
$\tau$, and local Haar measures multiply to the relevant global
Tamagawa measure. The section $\Phi_{s,v}$ is the holomorphic
Siegel--Weil section obtained from
$\rho_{1,3}(\varphi_{1,v}\otimes\bar\varphi_{2,v})$. With these
conventions,
\[
\begin{split}
 Z_v(1)&=\int_{T(F_v)}\Phi_{1,v}(\iota(t,1))
               \langle\tau_v(t)f_{1,v},f_{2,v}\rangle_v\,dt,\\
 Z_v^{\#}(1)&=
 \frac{L_v(3,\epsilon_{E_v/F_v})}
 {L_{E_v}(3/2,\operatorname{BC}(\tau_v)\otimes\gamma_v^3)}Z_v(1),\\
 \Phi_{1,v}(\iota(t,1))&=
       \langle\omega_v(t)\varphi_{1,v},\varphi_{2,v}\rangle_v.
\end{split}
\]
The normalized local factors equal one for normalized unramified data
outside a finite set. These conventions, including the last matrix
coefficient identity, are in Haan Sections 2.1, 2.4, and 3.1--3.2,
pp.3--8. The $L$-values at $3/2$ and $3$ are nonzero for the unitary
characters occurring here. Thus choosing nonzero local factors gives a
nonzero global lift. For the anisotropic rank-three group in the present
application, square integrability and cuspidality cause no additional
global convergence issue. Local occurrence and local test vectors still
must be checked; the formula does not assert that every prescribed
Schwartz vector has nonzero norm.

This displayed formula fixes one splitting convention completely.
To transport it to arbitrary splitting pairs, first compute the actual
unitary characters $\lambda(g)$ and $\kappa(t)$ in an identity
$\omega'(g,t)=\lambda(g)\kappa(t)\omega(g,t)$. Then the new lift with
input $\bar\tau$ is the old lift with input
$\overline{\tau\kappa^{-1}}$, multiplied by $\lambda$; its norm is
unchanged by the latter unitary twist. This gives the correct character
substitution without silently choosing a conjugation convention.
The factor $1/2$ in Haan's symplectic form must likewise be matched by
the form or additive-character normalization. Replacing both quotient
measures by probability measures multiplies the displayed target norm
by $\operatorname{vol}([G])^{-1}\operatorname{vol}([T])^{-2}$ when the
underlying input function is kept fixed. One must also account for
passage from $\U(V)$ and all adelic components to a connected
$\SU(V)$ quotient. No such factor is implicitly set equal to one here.

\subsection{THETA3-A1: the anisotropic Siegel zero-orbit identity}\label{A1}

\paragraph{Statement.}
Let $E/F$ be CM and let $V$ be an anisotropic Hermitian $E$-space.
Let $r\geq1$, and let $W^{\square}$ be split skew-Hermitian of rank
$2r$, with chosen complementary maximal isotropic subspaces. Use its
Schr\"odinger model $\mathcal S(V(\A_E)^r)$ and compatible splittings
for which
\[
 (\omega(g)\Phi)(x)=\Phi(g^{-1}x),\qquad
 (\omega(n(b))\Phi)(x)=\psi(\langle b,Q_V(x)\rangle)\Phi(x).
\]
Here $N$ is the Siegel unipotent, identified with the additive
$F$-space of Hermitian $r$ by $r$ matrices, $Q_V(x)$ is the Gram
matrix, and the pairing is nondegenerate. A conventional factor
$1/2$ may be absorbed in this identification. Give $[G]$ and $[N]$
probability Haar measures, where $G=\U(V)$, and define
\[
 T_\Phi(h)=\int_{[G]}\sum_{x\in V(E)^r}(\omega(g,h)\Phi)(x)\,dg.
\]
Then
\begin{equation}\label{A1-eq}
                    \int_{[N]}T_\Phi(n)\,dn=\Phi(0).
\end{equation}

\paragraph{Proof.}
Anisotropy makes $[G]$ compact; the centre is an anisotropic norm-one
torus as well. The quotient $[N]$ of an adelic vector group is compact.
On compact sets of $g$, the finite support condition confines rational
points to finitely many lattice cosets, and uniform archimedean
Schwartz decay dominates their polynomial counting function. Thus the
theta series converges absolutely and uniformly there. The $N$-action
adds only phases of modulus one. Summation and both integrals may
therefore be interchanged.

Since $Q_V(g^{-1}x)=Q_V(x)$, the inner unipotent integral is
\[
 \int_{[N]}\psi(\langle b,Q_V(x)\rangle)\,db
 =\begin{cases}1&Q_V(x)=0,\\0&Q_V(x)\ne0.\end{cases}
\]
For a nontrivial character the vanishing follows by translation by an
element on which its value differs from one. Nondegeneracy of the
pairing ensures that the character is trivial precisely for zero Gram
matrix. If the Gram matrix vanishes, every $x_i$ has zero norm, so
anisotropy forces $x=0$. The surviving term is fixed by $G$ and is
$\Phi(0)$, proving the formula. $\square$

The identity applies to both the ball Hermitian space and a totally
definite Hermitian space. It identifies the normalized zero-orbit
functional, not their nonzero Fourier coefficients, unaveraged periods,
or motives. A residual diagonal character changes its right-hand side
to $\Phi(0)\int_{[G]}\chi(g)\,dg$; isotropic $V$ instead admits
nonzero zero-Gram tuples. These hypotheses cannot be omitted.

\subsection{THETA3-A2: finite translation and quantitative nonvanishing}\label{A2}

\paragraph{Statement.}
In A1 write $\Phi=\Phi_\infty\otimes\Phi_f$, assume
$a=\Phi_\infty(0)\ne0$ and $\Phi_f(0)=1$, and keep the infinite input
fixed. For every $0<\varepsilon<1$ there is $h_f\in N(\A_{F,f})$
such that
\[
 \left|T_{\Phi_\infty\otimes\omega_f(h_f)\Phi_f}(1)\right|
                   \geq(1-\varepsilon)|a|.
\]
The finite Schwartz norm is unchanged. Every compact subgroup of
$G(\A_{F,f})$ fixing the total doubled input still fixes its translate.

\paragraph{Proof.}
The normalized mean in A1 is $a$. Compactness and continuity give
$n=(n_\infty,n_f)$ with $|T_\Phi(n)|\geq|a|$. Additive weak
approximation gives $\gamma_j\in N(F)$ with
$\gamma_{j,\infty}\longrightarrow n_\infty$. Automorphy and
continuity imply
\[
 T_\Phi(\gamma_{j,\infty},n_f)
 =T_\Phi(1,\gamma_{j,f}^{-1}n_f)
 \longrightarrow T_\Phi(n).
\]
Choose $j$ making the error at most $\varepsilon|a|$. Commutation of
the dual-pair actions identifies the value with the translated-input
integral. The finite Weil action is unitary. The same commutation shows
that every old $G$-stabilizer fixes the translated total input.
$\square$

This assertion changes finite theta data. Its finite tensor factors
may require a smaller common level after decomposition. It supplies
no uniform bound for their number or coefficients. The explicit
additive grid and Gaussian derivative certificate in THETA3-Q7
strengthen this existence argument under named effective data.

\subsection{THETA3-A3: the exact finite extraction cost}\label{A3}

\paragraph{Statement.}
Suppose an inverse-doubling decomposition and the finite torus
expansions of van Hoften--Sempliner Lemma 4.2.10 give
\[
 P=\sum_{a=1}^{R}c_a
       \sum_{\chi_i\in\mathcal X_{i,a},\ 1\leq i\leq4}
                       P_{a,\chi_1,\chi_2,\chi_3,\chi_4}.
\]
The inner terms are actual periods of integrated character lifts, with
conjugation on the last two factors and one fixed geometric
normalization. If $P\ne0$, at least one satisfies
\[
 |P_{a,\chi_1,\chi_2,\chi_3,\chi_4}|
 \geq \frac{|P|}{\sum_a|c_a|\prod_{i=1}^4|\mathcal X_{i,a}|}.
\]

\paragraph{Proof.}
Let $M$ be the maximum modulus of the finitely many periods in the
sum. The triangle inequality gives
$|P|\leq M\sum_a|c_a|\prod_i|\mathcal X_{i,a}|$; divide by the
positive denominator. $\square$

For an effective certificate use finite permissible character sets or
computable upper bounds for their sizes; determining the exact nonzero
support may itself require zero testing. Further rational Hodge-map or
connected-component extraction has its own finite coefficient cost.
The estimate is the same finite-sum principle developed in THETA3-Q5.

\subsection{THETA3-A4: higher-rank ball tensors and a surface section}\label{A4}

\paragraph{Statement.}
Replace the rank-three $V$ in the analytic part of 032 by a Hermitian
space of rank $m\geq3$, of signature $(m-1,1)$ at one real place and
positive definite at the others. Write $d=m-1$. For a multiplicity line
the first projected Taylor tensor again takes values in $\Omega_D^1$;
for a multiplicity plane, double alternation of the second tensor takes
values in $\Omega_D^2$. Finite actual inputs span every cotangent
fibre, and decomposable inputs satisfy the same wedge identity as in
032. Balanced compact multiplicity signs therefore give a nonzero
mixed pairing
\begin{equation}\label{A4-pair}
 \int_X\xi\wedge\bar\eta\wedge c_1(A)^{d-2}\ne0
\end{equation}
on a sufficiently deep compact principal ball quotient $X$, where
$A$ is ample and $\xi,\eta$ are the rational-decomposition endpoint
two-forms. A smooth complete-intersection surface $S\subset X$ of
divisors in $|A^{n_j}|$, $1\leq j\leq d-2$, satisfies
\[
 \int_S\xi|_S\wedge\overline{\eta|_S}
 =\left(\prod_jn_j\right)
      \int_X\xi\wedge\bar\eta\wedge c_1(A)^{d-2}\ne0.
\]
For $d=2$ take $S=X$ and the empty product one.

\paragraph{Proof.}
Let $L$ be the negative tautological line and $Q=V_\mu/L$, of rank
$m-1$. For multiplicity rank $k$ the active Lagrangian is
$L\otimes P\oplus Q^*\otimes P^*$. The determinant identity is
\[
 \det\mathcal F_P
 =L^{2k}\otimes(\det V_\mu)^{-k}\otimes(\det P)^{2-m}.
\]
The last two factors, and all compact-place factors, are constant on
$D$. On $\SU(V)$ the first of those constant representations is
trivial. After choosing the same constant-line trivializations as in
032, $(\det\mathcal F_P)^{1/2}=L^k$ times a constant line. Thus
\[
 (\det\mathcal F_{P_1})^{1/2}\otimes Q^*\otimes P_1^*
       \simeq L\otimes Q^*=\Omega_D^1
\]
for a line, and
\[
 (\det\mathcal F_P)^{1/2}\otimes
             \bigwedge^2Q^*\otimes\bigwedge^2P^*
       \simeq L^2\otimes\bigwedge^2Q^*=\Omega_D^2
\]
for a plane. The latter is the canonical bundle only when $m=3$.

The Gaussian-isolation argument of 032 Proposition 4.5 uses a position
space of dimension $m$ instead of three. In a fixed graph frame,
indicators of $x_0+k\widehat{\OO}_F^{\,m}$ give first tensors tending
to a nonzero Gaussian scalar times the Minkowski vector $x_0$.
Uniform Gaussian decay proves the convergence. These lattice vectors
span the complexified position space; the image of the span of actual
inputs is a finite-dimensional, hence closed, subspace. Projection
onto $L\otimes Q^*$ is surjective. Consequently actual theta values
span every cotangent fibre, and a finite selected list has a common
principal level. The covariance argument uses the same unitary Weil
splittings and triviality of determinant twists on $\SU(V)$, and is
independent of the value three.

For a multiplicity decomposition the completed theta series multiply.
The degree-two coefficient and double-alternation calculation stated
above is valid for any rank of $Q$. Select two first tensors with
independent values at one point; their wedge $\xi$ is nonzero.
Compactness of $X$ follows from a definite real place, and projectivity
from the arithmetic ball-quotient construction. Choose a positive
curvature representative $\omega$ of $c_1(A)$. Pointwise
$\xi\wedge\bar\xi\wedge\omega^{d-2}$ is a positive constant
times $\|\xi\|^2\omega^d$, so its integral is positive.

Keep the rational total multiplicity plane and its finite input fixed.
On a compact parameter neighborhood, the same Gaussian domination
with polynomial derivatives proves $C^\infty$ continuity of the
projected second tensor, exactly as in 032 Proposition 4.6(iii).
Therefore its mixed pairing with $\xi$, using $\omega^{d-2}$, remains
nonzero near the original decomposition. Balanced signs permit a
rational orthogonal decomposition with the requested endpoint signs
arbitrarily close to this compact-place datum, by the same weak
approximation argument. Only multiplicity subspaces move; $V$ is
unchanged. At that rational endpoint the tensor is a finite sum of
wedges of holomorphic first tensors, so is holomorphic and closed.
Pass to a common principal cover for the selected finite families.
The integral scales by the positive covering degree. This proves
\eqref{A4-pair}.

For sufficiently large $n_j$, Bertini gives a smooth projective
complete-intersection surface. Its cohomology class is
$(\prod_jn_j)c_1(A)^{d-2}$. Both endpoint forms are closed because
holomorphic forms on a projective manifold are closed. The
cohomological restriction formula gives the displayed surface
identity. This also shows directly that every such smooth complete
intersection retains a nonzero holomorphic two-form when $\eta=\xi$.
$\square$

This extends the analytic and geometric wedge mechanism. It does not
by itself extend 032's precise $U(v)$ arithmetic extraction theorem
with its rank-three Frobenius calculation. The exact all-rank source
identification is supplied separately in A6 using Liu's componentwise
Theorems 4.15 and 4.18 and the reflex-type calculation. Liu's
Corollary 4.20 also gives the CM Albanese conclusion in
all ranks $m\geq3$, so simply raising this ball rank does not create
non-CM $H^1$. Nor does adding positive directions turn this family
into a hyperbolic Witt tower or a chain of totally definite theta
lifts. A theorem imposing signatures at every real place still
requires those signatures to match: one active ball and positive
definite compact places do not satisfy such a hypothesis merely by
increasing $m$.

\subsection{THETA3-A5: an exact finite sum of definite Bessel periods}\label{A5}

\paragraph{Setup.}
Return to rank three and let $G=\U(V)$, with the ball signature at one
real place and positive definite signature elsewhere. Choose
$z\in V(E)$ negative at the active place, using weak approximation;
then $H=\U(z^\perp)$ is totally definite of rank two. View $H$ as a
subgroup of $G$ fixing $z$. Let $K_f\subset G(\A_{F,f})$ be compact
open and let $K_\infty$ be the stabilizer of the negative line $Ez$ at
the active place, times the compact groups at the other real places.
Fix a unitary scalar character $\kappa$ of $K_\infty$.

Let $\mathcal E$ be a nonzero finite-dimensional space of smooth
automorphic functions on $[G]$, right $K_f$-invariant, satisfying
$f(gk)=\kappa(k)^{-1}f(g)$ for $k\in K_\infty$. Give it the
Petersson inner product inherited from probability measure on $[G]$.
The application is a fixed-level span of holomorphic canonical theta
forms in their equivariant adelic realization. On a ball surface the
canonical fibre is one-dimensional, which gives precisely a scalar
$K_\infty$ type. Fix the splitting and central type and use the full
adelic realization, as in van Hoften--Sempliner, before extracting
connected components. This is not an assertion that an arbitrary
unrealized form on a connected $\SU(V)$ quotient has already been
extended to $G$.

For $a\in G(\A_{F,f})$ put
\[
 R_af(h)=f(ha),\qquad
 K_H(a)=H(\A_{F,f})\cap aK_fa^{-1}.
\]
Give $[H]$ probability measure, and let $\mathcal M_a$ be the space of
automorphic functions on $H$ right $K_H(a)$-invariant and of right
$H(F\otimes_\Q\R)$-type $\kappa^{-1}|_H$.

\paragraph{Statement.}
Writing $D=\dim_\C\mathcal E$, there are $s\leq D$ finite adelic
elements $a_1,\ldots,a_s$ for which
\[
 \mathcal Q=\sum_{j=1}^sR_{a_j}^*R_{a_j}
\]
is positive definite on $\mathcal E$. For all $u,v\in\mathcal E$,
\begin{equation}\label{A5-frame}
 \langle u,v\rangle_G
  =\sum_{j=1}^s
          \langle R_{a_j}\mathcal Q^{-1}u,R_{a_j}v\rangle_H.
\end{equation}
Choose an orthonormal basis $e_{j,\alpha}$ of $\mathcal M_{a_j}$ and
put
\[
 \mathcal P_{j,\alpha}(f)=
    \int_{[H]}f(ha_j)\overline{e_{j,\alpha}(h)}\,dh.
\]
With inner products conjugate-linear in the first variable, the
identity becomes the finite sum
\begin{equation}\label{A5-periods}
 \langle u,v\rangle_G
  =\sum_{j,\alpha}
       \overline{\mathcal P_{j,\alpha}(\mathcal Q^{-1}u)}
                   \mathcal P_{j,\alpha}(v).
\end{equation}
These are genuine global adjacent-unitary-group periods of $\U(3)$
against totally definite $\U(2)$. Moreover
\begin{equation}\label{A5-bounds}
 \lambda_{\min}(\mathcal Q)\|u\|_G^2
 \leq\sum_{j,\alpha}|\mathcal P_{j,\alpha}(u)|^2
 \leq\lambda_{\max}(\mathcal Q)\|u\|_G^2.
\end{equation}

\paragraph{Proof: well-defined finite restriction spaces.}
The left $H(F)$-invariance is inherited from $G(F)$-invariance. If
$k\in K_H(a)$, then $a^{-1}ka\in K_f$, so
$f(hka)=f(ha)$. At infinity, finite $a$ commutes with $H(F\otimes\R)$,
which is contained in $K_\infty$; this proves the asserted right type.
Thus $R_a$ maps into $\mathcal M_a$.

Because $H$ is totally definite,
$H(F)\backslash H(\A_{F,f})/K_H(a)$ is finite. After choosing finite
representatives, a scalar-type function is determined by its values
at those representatives. A finite stabilizer can force a value to
vanish when its character is nontrivial, but introduces no infinite
dimensional space. Thus $\mathcal M_a$ is finite-dimensional. Its
norm is a finite weighted sum of squared values, with positive masses
proportional to the reciprocal orders of the finite stabilizers. In
particular, a smooth nonzero restriction has positive norm.

\paragraph{Proof: separation and at most $D$ translates.}
The group $G$ satisfies weak approximation at the archimedean places
needed here. A direct proof uses the Cayley transform
\[
 X\longmapsto(1-X)(1+X)^{-1},\qquad X^*=-X.
\]
The skew-adjoint matrices form an $F$-vector space; its $F$-points
are dense in the product of its real completions. The Cayley transform
is an isomorphism of open sets, with inverse
$(1-g)(1+g)^{-1}$. Its real image, excluding the proper locus with
eigenvalue $-1$, is dense. Hence $G(F)$ is dense in
$G(F\otimes_\Q\R)$. This uses $\U(V)$, not an unproved
approximation assertion for a chosen $\SU(V)$ presentation.

If $f\in\mathcal E$ is nonzero, choose $(g_\infty,g_f)$ where its
value is nonzero. Continuity and weak approximation provide
$\gamma\in G(F)$ such that $f(\gamma_\infty,g_f)\ne0$.
Automorphy gives
\[
 f(\gamma_\infty,g_f)=f(1,\gamma_f^{-1}g_f).
\]
For $a=\gamma_f^{-1}g_f$, therefore $R_af(1)\ne0$, and
$R_af$ has positive norm. The family of restrictions separates
$\mathcal E$.

Start with its full kernel space $\mathcal E_0=\mathcal E$.
If $\mathcal E_i\ne0$, choose nonzero $f_i\in\mathcal E_i$ and
one restriction not annihilating it, and intersect $\mathcal E_i$
with that restriction's kernel. Dimension drops strictly. After at
most $D$ steps the common kernel is zero. Thus
$\langle u,\mathcal Q u\rangle_G=\sum_j\|R_{a_j}u\|_H^2$
is positive for every nonzero $u$.

\paragraph{Proof: identity, Parseval, and bounds.}
The operator $\mathcal Q$ is positive and self-adjoint, hence
invertible. By its definition,
\[
 \sum_j\langle R_{a_j}\mathcal Q^{-1}u,R_{a_j}v\rangle_H
 =\langle\mathcal Q^{-1}u,\mathcal Q v\rangle_G
 =\langle u,v\rangle_G.
\]
Finite-dimensional Parseval in each $\mathcal M_{a_j}$ gives
\eqref{A5-periods}. The spectral bounds for a positive self-adjoint
operator give \eqref{A5-bounds}. $\square$

\paragraph{Exact matrices and conditional quantitative bounds.}
In a basis of $\mathcal E$, let $G_{\rm ball}$ be its Petersson Gram
matrix. Let $R_j$ be the finite evaluation matrix and $D_j$ the
positive diagonal mass matrix for $\mathcal M_{a_j}$. Then
\[
 H_{\rm fin}=\sum_jR_j^*D_jR_j,\qquad
                  \mathcal Q=G_{\rm ball}^{-1}H_{\rm fin}.
\]
This formula keeps the analytic Gram matrix separate from definite
finite masses. If $D>1$, $G_{\rm ball}\leq M I$, and a square
$D$ by $D$ minor $B$ of the stacked matrix with blocks
$D_j^{1/2}R_j$ has determinant of modulus at least $\delta>0$, then
Cauchy--Binet and the eigenvalue product give the valid lower bound
\[
 \lambda_{\min}(\mathcal Q)
 \geq\frac{\delta^2}{M\bigl(\operatorname{tr}H_{\rm fin}\bigr)^{D-1}}.
\]
The case $D=1$ is immediate. A height or denominator argument can
supply $\delta$ only after algebraicity of appropriately normalized
entries has actually been proved. No such algebraicity, and no
rationality of $\mathcal Q^{-1}$, follows from finite dimensionality.

\paragraph{Meaning of the comparison.}
Formula \eqref{A5-periods} is an exact finite-sum refinement of the
desired comparison with definite periods, with an explicitly defined
frame correction. A fixed restriction alone computes its own norm,
not the ball norm. The evaluation of canonical forms at these CM
points is a fibrewise operation, not a pullback of a degree-two
cohomology class to a zero-dimensional variety. The analytic operator
$\mathcal Q^{-1}$ is not automatically an algebraic correspondence.
An algebraic, rational, or uniform-integral realization is an
additional theorem required for a motivic transfer.

\subsection{Direct sums, genuine seesaws, and iterated lifts}

For compatible splittings and $W=W_1\perp W_2$, the common oscillator
representation gives
\[
 \theta_W(g,h_1\oplus h_2;\varphi_1\otimes\varphi_2)
 =\theta_{W_1}(g,h_1;\varphi_1)
             \theta_{W_2}(g,h_2;\varphi_2).
\]
The differential version is the double-alternating wedge identity.
It belongs to the actual seesaw whose pairs are
$\U(V)\times\U(V)$ against $\U(W_1)\times\U(W_2)$, and diagonal
$\U(V)$ against $\U(W)$. Integrating against characters or test
automorphic forms gives a seesaw identity by Fubini on the compact
quotients in this application. A second orthogonal decomposition of
the same $W$ supplies two other tori in the same group. It does not
change $V$ to a definite Hermitian space.

This construction is bilinear in two independent inputs and combines
two kernels by a direct sum. An iterated lift is instead a composition
of linear integral operators, with the output automorphic
representation of one step serving as the input to the next. One
cannot identify the two merely from their number of theta factors.
The literal step 8 of Alp\"oge's note, p.4, describes a chain of
totally definite groups and changes of automorphic parameters. Its
recovered mirror contains actual successive operators \texttt{up}
and \texttt{dn}, rather than the above wedge operation. A comparison
therefore needs intertwining maps and matching representations,
splittings, and vectors. The finite-frame identity A5 and the
restriction seesaw for $V=z^\perp\perp Ez$ are concrete alternative
comparisons; neither asserts that the two independent degree-one
lifts already form that iterated chain.
